\section{Conclusions and Further Scope for Research}
By extending mean-field signal propagation to include multiplicative dropout, we established that activation non-analyticity fundamentally dictates universality. This formalism yields a Landau equation of state for the decorrelation order parameter, a universal two-parameter scaling collapse, and distinct critical exponents for smooth versus kinked channels. Mean-field theory is often best at clarifying the mechanisms inside wide networks, and less often at producing directly testable design rules; here it motivates one such rule. Treating dropout as a depth-dependent field gives a local-stationary propagation objective favoring saturated allocation, and an additive regularization-reach heuristic places that allocation near the input. Their use for abrupt schedules and trained networks is tested empirically. These schedules reduce final test loss by \(18\)--\(35\%\) in MLP settings and by \(4\)--\(6\%\) in Vision Transformer settings relative to constant dropout at the same budget (Table~\ref{tab:loss_improvements}), with secondary accuracy gains and no extra computational cost. Comparing minimum recorded test loss gives smaller reductions of \(0.9\)--\(2.5\%\) for the same schedules. The growth of the advantage during continued training is itself a striking result: in the CIFAR-10 budget comparison, front-loading substantially restrains the late rise in test loss while accuracy continues to improve, suggesting that its practical benefit includes limiting overfitting after the loss minimum. The broader comparisons select checkpoints and profiles from validation loss and show smaller, task-dependent gains, including losses to uniform dropout in the extended Jannis Transformer cohort. Extending this mean-field analysis to attention mechanisms and incorporating finite-width corrections are natural directions for future work.

\section{Limitations}
Our analysis is the leading-order, large-width description: finite-width corrections, controlled perturbatively by \(L/N\), are not computed explicitly. They can accumulate with depth; we have not established whether they preserve the leading exponents or how they change a useful scheduling profile. The scheduling objective also freezes local coefficients and replaces trajectory-dependent derivatives by stationary values, while the reach argument assumes additive exposure at a common decay scale. Finite dropout, multiple masks and feature learning can change both approximations.

Architectural scope is also restricted. We treat dropout most explicitly for MLPs; App.~\ref{app:cnn_extension} argues that infinite-channel CNNs and residual branches inherit the same smooth/kinked split through the local Gaussian kernel, but we do not derive the full spatial-mode CNN Landau theory, run CNN experiments, or develop a dropout-deformed ResNet theory. Skip connections in particular modify large-depth propagation while preserving the local channel, so promoting residual-branch strength to a control parameter would put the Transformer experiments on firmer footing.

Finally, the theory is an initialization theory of \emph{forward} correlation dynamics. We give the leading backward covariance recursion and note that dropout creates the same diagonal/off-diagonal asymmetry for gradients, but do not develop the full backward critical theory (finite-width gradient susceptibilities, training-time mask correlations, feature-learning effects), nor do we model how observables evolve after feature learning substantially reshapes representations, e.g.\ in catapult regimes \cite{zhu2024quadraticcatapult}. Extending the same approach to training-time regularizers, dropout warm-up, adaptive schedules, \(L^2\) penalties, or broader scaling-law questions \cite{Bahri2024ExplainingScalingLaws} is a natural direction for future work.

\noindent\emph{Acknowledgments.} The author thanks Riccardo Penco for support and discussions, and the anonymous ICML reviewers for feedback that sharpened the scheduling discussion and the treatment of the kinked-class scaling path. This work is partially supported by the U.S. Department of Energy under grant DE-SC0010118.

\section*{Reproducibility Statement}
The repository linked in the introduction contains the original notebooks and saved results, the manuscript source, and the paired seed endpoints and analysis script underlying Tables~\ref{tab:loss_improvements} and~\ref{tab:frontloaded_results}. App.~\ref{app:experimental_info} reports the original run counts, hyperparameters, dataset splits, hardware, and sweep details. App.~\ref{app:extended_experiments} reports the additional benchmark and profile-pilot evidence, including incomplete coverage. The \texttt{benchmarks/} directory contains the retained benchmark training and data-preparation code, frozen profile configurations and split identifiers; App.~\ref{app:benchmark_protocols} documents their protocols and the unrecovered settings in the standalone financial study. The two main gain tables report nominal $95\%$ confidence intervals; existing curve bands and $\pm$ summaries retain their SEM convention. The mean-field fits use deterministic recursion code from the same repository and report fit standard errors separately. The CIFAR-100 archive lacks the configuration needed to verify all reported settings (Table~\ref{tab:vit_config}).

\section*{Impact Statement}
The main practical benefit is improved performance at fixed training cost, or potentially reduced training cost when targeting a fixed performance level. We do not anticipate additional ethical risks.
